\documentclass[aspectratio=169]{beamer}
\usepackage{../pomozno/prosojnice}
\usepackage{fontawesome5}

% Izpis dokumenta na prosojnicah. Številke so natipkane, da kažejo, kaj bi
% izpisal LaTeX, tudi kadar je izpis napačen.
\newcommand{\razdelek}[2]{\par\smallskip{\large\bfseries #1\quad #2}\par\smallskip}
\newcommand{\izrek}[2]{\textbf{Izrek #1 (#2).}}
\newcommand{\izrekams}[2]{\textbf{Izrek #1} (#2)\textbf{.}}

\begin{document}
\begin{frame}
  \begin{center}
    \Huge\textbf{Sklicevanje in okolja}
  \end{center}
\end{frame}

\section{Sklici}

\begin{frame}[fragile]{}
  \begin{columns}[T]
    \begin{column}{0.55\textwidth}
      \begin{minted}[fontsize=\scriptsize]{latex}
\section{Uvod}
Trikotnike poznamo že iz osnovne šole.

\section{Pitagorov izrek}
\textbf{Izrek 2.1 (Pitagorov izrek).}
\emph{V pravokotnem trikotniku s katetama
\(a\) in \(b\) ter hipotenuzo \(c\) velja
\(a^2 + b^2 = c^2\).}

Iz izreka 2.1 izpeljemo razdaljo med
točkama v prostoru~[1].

\section*{Literatura}
[1] I. Vidav, \emph{Višja matematika I}, 2008.
      \end{minted}
    \end{column}
    \begin{column}{0.4\textwidth}
      \footnotesize
      \razdelek{1}{Uvod}
      Trikotnike poznamo že iz osnovne šole.
      \razdelek{2}{Pitagorov izrek}
      \izrek{2.1}{Pitagorov izrek}
      \emph{V pravokotnem trikotniku s katetama \(a\) in \(b\) ter hipotenuzo \(c\) velja \(a^2 + b^2 = c^2\).}

      Iz izreka 2.1 izpeljemo razdaljo med točkama v prostoru~[1].
      \razdelek{}{Literatura}
      [1] I. Vidav, \emph{Višja matematika I}, 2008.
    \end{column}
  \end{columns}
\end{frame}

\begin{frame}[fragile]{}
  \begin{columns}[T]
    \begin{column}{0.55\textwidth}
      \begin{minted}[fontsize=\scriptsize,highlightlines={4-6}]{latex}
\section{Uvod}
Trikotnike poznamo že iz osnovne šole.

\section{Talesov izrek}
\textbf{Izrek 2.1 (Talesov izrek).}
\emph{Obodni kot nad premerom je pravi.}

\section{Pitagorov izrek}
\textbf{Izrek 2.1 (Pitagorov izrek).}
\emph{V pravokotnem trikotniku s katetama
\(a\) in \(b\) ter hipotenuzo \(c\) velja
\(a^2 + b^2 = c^2\).}
      \end{minted}
    \end{column}
    \begin{column}{0.4\textwidth}
      \footnotesize
      \razdelek{1}{Uvod}
      Trikotnike poznamo že iz osnovne šole.
      \razdelek{2}{Talesov izrek}
      \izrek{2.1}{Talesov izrek}
      \emph{Obodni kot nad premerom je pravi.}
      \razdelek{3}{Pitagorov izrek}
      \alert{\izrek{2.1}{Pitagorov izrek}}
      \emph{V pravokotnem trikotniku s katetama \(a\) in \(b\) ter hipotenuzo \(c\) velja \(a^2 + b^2 = c^2\).}
    \end{column}
  \end{columns}
\end{frame}

\begin{frame}[fragile]{}
  \begin{columns}[T]
    \begin{column}{0.55\textwidth}
      \begin{minted}[fontsize=\scriptsize]{latex}
\usepackage{amsthm}
\newtheorem{izrek}{Izrek}[section]
...
\begin{document}
...
\section{Pitagorov izrek}
\begin{izrek}[Pitagorov izrek]
  V pravokotnem trikotniku ...
\end{izrek}
      \end{minted}
    \end{column}
    \begin{column}{0.4\textwidth}
      \footnotesize
      \razdelek{2}{Pitagorov izrek}
      \izrekams{2.1}{Pitagorov izrek}
      \emph{V pravokotnem trikotniku \dots}
    \end{column}
  \end{columns}
\end{frame}

\begin{frame}[fragile]{}
  \begin{columns}[T]
    \begin{column}{0.55\textwidth}
      \begin{minted}[fontsize=\scriptsize,highlightlines={6-9}]{latex}
\usepackage{amsthm}
\newtheorem{izrek}{Izrek}[section]
...
\begin{document}
...
\section{Talesov izrek}
\begin{izrek}[Talesov izrek]
  Obodni kot nad premerom je pravi.
\end{izrek}

\section{Pitagorov izrek}
\begin{izrek}[Pitagorov izrek]
  V pravokotnem trikotniku ...
\end{izrek}
      \end{minted}
    \end{column}
    \begin{column}{0.4\textwidth}
      \footnotesize
      \razdelek{2}{Talesov izrek}
      \izrekams{2.1}{Talesov izrek}
      \emph{Obodni kot nad premerom je pravi.}
      \razdelek{3}{Pitagorov izrek}
      \izrekams{3.1}{Pitagorov izrek}
      \emph{V pravokotnem trikotniku \dots}
    \end{column}
  \end{columns}
\end{frame}

\begin{frame}[fragile]{}
  \begin{tabular}{ll}
    \mintinline{latex}|\newtheorem{izrek}{Izrek}| & Izrek 1, Izrek 2, Izrek 3 \\
    \mintinline{latex}|\newtheorem{izrek}{Izrek}[section]| & Izrek 2.1, Izrek 3.1, Izrek 3.2 \\
    \mintinline{latex}|\newtheorem{posledica}[izrek]{Posledica}| & Izrek 3.1, Posledica 3.2, Izrek 3.3 \\
  \end{tabular}
\end{frame}

\begin{frame}[fragile]{}
  \begin{columns}[T]
    \begin{column}{0.5\textwidth}
      \begin{minted}[fontsize=\scriptsize]{latex}
{\theoremstyle{plain}
\newtheorem{izrek}{Izrek}[section]
\newtheorem{posledica}[izrek]{Posledica}
}

{\theoremstyle{definition}
\newtheorem{definicija}[izrek]{Definicija}
}
      \end{minted}
    \end{column}
    \begin{column}{0.45\textwidth}
      \footnotesize
      \razdelek{3}{Pitagorov izrek}
      \textbf{Definicija 3.1.}
      Trikotnik je \emph{pravokoten}, če je eden od njegovih kotov pravi.
      Stranici ob pravem kotu sta kateti, tretja stranica je hipotenuza.

      \medskip
      \izrekams{3.2}{Pitagorov izrek}
      \emph{V pravokotnem trikotniku s katetama \(a\) in \(b\) ter hipotenuzo \(c\) velja \(a^2 + b^2 = c^2\).}

      \medskip
      \textbf{Posledica 3.3.}
      \emph{V pravokotnem trikotniku je hipotenuza daljša od vsake od katet.}
    \end{column}
  \end{columns}
\end{frame}

\begin{frame}[fragile]
  \begin{columns}[T]
    \begin{column}{0.55\textwidth}
      \begin{minted}[fontsize=\scriptsize]{latex}
\section{Pitagorov izrek}
\begin{izrek}[Pitagorov izrek]
  V pravokotnem trikotniku ...
\end{izrek}

Iz izreka~2.1 izpeljemo razdaljo med
točkama v prostoru~[1].
      \end{minted}
    \end{column}
    \begin{column}{0.4\textwidth}
      \footnotesize
      \razdelek{2}{Pitagorov izrek}
      \izrekams{2.1}{Pitagorov izrek}
      \emph{V pravokotnem trikotniku \dots}

      Iz izreka~2.1 izpeljemo razdaljo med točkama v prostoru~[1].
    \end{column}
  \end{columns}
\end{frame}

\begin{frame}[fragile]
  \begin{columns}[T]
    \begin{column}{0.55\textwidth}
      \begin{minted}[fontsize=\scriptsize,highlightlines={1-4}]{latex}
\section{Talesov izrek}
\begin{izrek}[Talesov izrek]
  Obodni kot nad premerom je pravi.
\end{izrek}

\section{Pitagorov izrek}
\begin{izrek}[Pitagorov izrek]
  V pravokotnem trikotniku ...
\end{izrek}

Iz izreka~2.1 izpeljemo razdaljo med
točkama v prostoru~[1].
      \end{minted}
    \end{column}
    \begin{column}{0.4\textwidth}
      \footnotesize
      \razdelek{2}{Talesov izrek}
      \izrekams{2.1}{Talesov izrek}
      \emph{Obodni kot nad premerom je pravi.}
      \razdelek{3}{Pitagorov izrek}
      \izrekams{3.1}{Pitagorov izrek}
      \emph{V pravokotnem trikotniku \dots}

      Iz izreka~\alert{2.1} izpeljemo razdaljo med točkama v prostoru~[1].
    \end{column}
  \end{columns}
\end{frame}

\begin{frame}[fragile]{}
  \begin{columns}[T]
    \begin{column}{0.55\textwidth}
      \begin{minted}[fontsize=\scriptsize,highlightlines={3,8}]{latex}
\section{Pitagorov izrek}
\begin{izrek}[Pitagorov izrek]
  \label{izr:pitagora}
  V pravokotnem trikotniku s katetama \(a\)
  in \(b\) ter hipotenuzo \(c\) velja ...
\end{izrek}

Iz izreka~\ref{izr:pitagora} izpeljemo
razdaljo med točkama v prostoru~[1].
      \end{minted}
    \end{column}
    \begin{column}{0.4\textwidth}
      \footnotesize
      \razdelek{3}{Pitagorov izrek}
      \izrekams{3.1}{Pitagorov izrek}
      \emph{V pravokotnem trikotniku s katetama \(a\) in \(b\) ter hipotenuzo \(c\) velja \dots}

      Iz izreka 3.1 izpeljemo razdaljo med točkama v prostoru~[1].
    \end{column}
  \end{columns}
  \bigskip
\end{frame}

\begin{frame}[fragile]
  \begin{center}
    \Large
    \mintinline{latex}|Iz izreka 2.1|
    \bigskip

    \mintinline{latex}|Iz izreka~\ref{izr:pitagora}|
  \end{center}
  \bigskip
\end{frame}

\begin{frame}[fragile]{}
  \begin{columns}[T]
    \begin{column}{0.55\textwidth}
      \begin{minted}{latex}
Iz izreka~\ref{izr:pitagor}
izpeljemo razdaljo med
točkama v prostoru~[1].
      \end{minted}
    \end{column}
    \begin{column}{0.4\textwidth}
      Iz izreka \textbf{??} izpeljemo razdaljo med točkama v prostoru~[1].
    \end{column}
  \end{columns}
  \bigskip
  \begin{minted}[fontsize=\footnotesize]{text}
LaTeX Warning: Reference `izr:pitagor' on page 2 undefined on input line 31.
  \end{minted}
\end{frame}

\begin{frame}[fragile]{}
  \begin{minted}{bash}
$ pdflatex dokument.tex
$ pdflatex dokument.tex
  \end{minted}
\end{frame}

\begin{frame}[fragile]{}
  \begin{columns}[T]
    \begin{column}{0.45\textwidth}
      \begin{minted}{latex}
\[
  a^2 + b^2 = c^2
\]
      \end{minted}
      \[
        a^2 + b^2 = c^2
      \]
    \end{column}
    \begin{column}{0.45\textwidth}
      \begin{minted}{latex}
\begin{equation}
  a^2 + b^2 = c^2
\end{equation}
      \end{minted}
      \begin{equation}
        a^2 + b^2 = c^2
      \end{equation}
    \end{column}
  \end{columns}
\end{frame}

\begin{frame}[fragile]{}
  \begin{minted}{latex}
\begin{equation}
  \label{enc:pitagora}
  a^2 + b^2 = c^2
\end{equation}
  \end{minted}
  \medskip

  \begin{tabular}{ll}
    \mintinline{latex}|izreka~\ref{izr:pitagora}| & izreka 3.1 \\
    \mintinline{latex}|enačbe~\eqref{enc:pitagora}| & enačbe (1) \\
    \mintinline{latex}|na strani~\pageref{izr:pitagora}| & na strani 2 \\
  \end{tabular}
  \medskip

  Predpone \texttt{izr:}, \texttt{enc:}, \texttt{tab:} so dogovor.
\end{frame}

\begin{frame}[fragile]{}
  \begin{align*}
    a x^2 + b x + c &= 0 \\
    \left(x + \frac{b}{2a}\right)^2 &= \frac{b^2-4ac}{4a^2} \\
    x_{1, 2} &= \frac{-b \pm \sqrt{b^2-4ac}}{2a}
  \end{align*}

  \begin{minted}{latex}
\begin{align*}
  a x^2 + b x + c &= 0 \\
  \left(x + \frac{b}{2a}\right)^2 &= \frac{b^2-4ac}{4a^2} \\
  x_{1, 2} &= \frac{-b \pm \sqrt{b^2-4ac}}{2a}
\end{align*}
  \end{minted}
\end{frame}

\begin{frame}[fragile]{Če nočete oštevilčenja, uporabite \texttt{align*}}
  \begin{dobroenv}
    \begin{minted}[fontsize=\scriptsize]{latex}
\begin{align*}
  a_{11} x_1 + a_{12} x_2 + \dots + a_{1n} x_n &= b_1 \\
  a_{21} x_1 + a_{22} x_2 + \dots + a_{2n} x_n &= b_2 \\
                                               &\vdotswithin{=} \\
  a_{n1} x_1 + a_{n2} x_2 + \dots + a_{nn} x_n &= b_n,
\end{align*}
    \end{minted}
    \small
    \begin{align*}
      a_{11} x_1 + a_{12} x_2 + \dots + a_{1n} x_n &= b_1 \\
      a_{21} x_1 + a_{22} x_2 + \dots + a_{2n} x_n &= b_2 \\
      &\vdotswithin{=} \\
      a_{n1} x_1 + a_{n2} x_2 + \dots + a_{nn} x_n &= b_n,
    \end{align*}
  \end{dobroenv}
\end{frame}

\begin{frame}[fragile]{Če nočete oštevilčenja, uporabite \texttt{align*}}
  \begin{slaboenv}
    \begin{minted}[fontsize=\scriptsize]{latex}
\begin{align}
  a_{11} x_1 + a_{12} x_2 + \dots + a_{1n} x_n &= b_1 \\
  a_{21} x_1 + a_{22} x_2 + \dots + a_{2n} x_n &= b_2 \\
                                               &\vdotswithin{=} \\
  a_{n1} x_1 + a_{n2} x_2 + \dots + a_{nn} x_n &= b_n,
\end{align}
    \end{minted}
    \small
    \begin{align}
      a_{11} x_1 + a_{12} x_2 + \dots + a_{1n} x_n &= b_1 \\
      a_{21} x_1 + a_{22} x_2 + \dots + a_{2n} x_n &= b_2 \\
      &\vdotswithin{=} \\
      a_{n1} x_1 + a_{n2} x_2 + \dots + a_{nn} x_n &= b_n,
    \end{align}
  \end{slaboenv}
\end{frame}

\begin{frame}[fragile]{Če nočete oštevilčenja, uporabite \texttt{align*}}
  \begin{slaboenv}
    \begin{minted}[fontsize=\scriptsize]{latex}
\begin{align}
  a_{11} x_1 + a_{12} x_2 + \dots + a_{1n} x_n &= b_1 \nonumber \\
  a_{21} x_1 + a_{22} x_2 + \dots + a_{2n} x_n &= b_2 \nonumber \\
                                               &\vdotswithin{=} \nonumber \\
  a_{n1} x_1 + a_{n2} x_2 + \dots + a_{nn} x_n &= b_n, \nonumber
\end{align}
    \end{minted}
    \small
    \begin{align}
      a_{11} x_1 + a_{12} x_2 + \dots + a_{1n} x_n &= b_1 \nonumber \\
      a_{21} x_1 + a_{22} x_2 + \dots + a_{2n} x_n &= b_2 \nonumber \\
      &\vdotswithin{=} \nonumber \\
      a_{n1} x_1 + a_{n2} x_2 + \dots + a_{nn} x_n &= b_n, \nonumber
    \end{align}
  \end{slaboenv}
\end{frame}

\begin{frame}[fragile]{Če uporabljate \texttt{align}, uporabite tudi \texttt{\&}}
  \begin{dobroenv}
    \begin{minted}[fontsize=\scriptsize]{latex}
\begin{align*}
  a_{11} x_1 + a_{12} x_2 + \dots + a_{1n} x_n &= b_1 \\
  a_{21} x_1 + a_{22} x_2 + \dots + a_{2n} x_n &= b_2 \\
                                               &\vdotswithin{=} \\
  a_{n1} x_1 + a_{n2} x_2 + \dots + a_{nn} x_n &= b_n,
\end{align*}
    \end{minted}
    \small
    \begin{align*}
      a_{11} x_1 + a_{12} x_2 + \dots + a_{1n} x_n &= b_1 \\
      a_{21} x_1 + a_{22} x_2 + \dots + a_{2n} x_n &= b_2 \\
      &\vdotswithin{=} \\
      a_{n1} x_1 + a_{n2} x_2 + \dots + a_{nn} x_n &= b_n,
    \end{align*}
  \end{dobroenv}
\end{frame}

\begin{frame}[fragile]{Če uporabljate \texttt{align}, uporabite tudi \texttt{\&}}
  \begin{slaboenv}
    \begin{minted}[fontsize=\scriptsize]{latex}
\begin{align*}
  a_{11} x_1 + a_{12} x_2 + \dots + a_{1n} x_n = b_1 \\
  a_{21} x_1 + a_{22} x_2 + \dots + a_{2n} x_n = b_2 \\
                                               \vdotswithin{=} \\
  a_{n1} x_1 + a_{n2} x_2 + \dots + a_{nn} x_n = b_n,
\end{align*}
    \end{minted}
    \small
    \begin{align*}
      a_{11} x_1 + a_{12} x_2 + \dots + a_{1n} x_n = b_1 \\
      a_{21} x_1 + a_{22} x_2 + \dots + a_{2n} x_n = b_2 \\
      \vdotswithin{=} \\
      a_{n1} x_1 + a_{n2} x_2 + \dots + a_{nn} x_n = b_n,
    \end{align*}
  \end{slaboenv}
\end{frame}

\begin{frame}[fragile]{Če uporabljate \texttt{align}, uporabite tudi \texttt{\&}}
  \begin{slaboenv}
    \begin{minted}[fontsize=\scriptsize]{latex}
\begin{align*}
  a_{11} x_1 + a_{12} x_2 + \dots + a_{1n} x_n = b_1 \\
  a_{21} x_1 + a_{22} x_2 + \dots + a_{2n} x_n = b_1 + b_2 \\
                                               \vdotswithin{=} \\
  a_{n1} x_1 + a_{n2} x_2 + \dots + a_{nn} x_n = b_1 + b_2 + \dots + b_n,
\end{align*}
    \end{minted}
    \small
    \begin{align*}
      a_{11} x_1 + a_{12} x_2 + \dots + a_{1n} x_n = b_1 \\
      a_{21} x_1 + a_{22} x_2 + \dots + a_{2n} x_n = b_1 + b_2 \\
      \vdotswithin{=} \\
      a_{n1} x_1 + a_{n2} x_2 + \dots + a_{nn} x_n = b_1 + b_2 + \dots + b_n,
    \end{align*}
  \end{slaboenv}
\end{frame}

\begin{frame}[fragile]{Tabele se po lahko premikajo po straneh}

  \begin{slaboenv}
    \begin{minted}{latex}
\begin{table}[!ht]
  \centering
  \caption{Pitagorejske trojice}
  \begin{tabular}{ccc}
    3 & 4 & 5 \\
    5 & 12 & 13 \\
  \end{tabular}
\end{table}

Prvih nekaj pitagorejskih trojic je v zgornji tabeli.
    \end{minted}
  \end{slaboenv}

\end{frame}

\begin{frame}[fragile]{Tabele se po lahko premikajo po straneh}

  \begin{dobroenv}
    \begin{minted}[highlightlines=4]{latex}
\begin{table}[!ht]
  \centering
  \caption{Pitagorejske trojice}
  \label{tab:trojice}
  \begin{tabular}{ccc}
    3 & 4 & 5 \\
    5 & 12 & 13 \\
  \end{tabular}
\end{table}

Prvih nekaj pitagorejskih trojic je v tabeli~\ref{tab:trojice}.
    \end{minted}
  \end{dobroenv}

\end{frame}

\section{Literatura}


\begin{frame}[fragile]{}
  \begin{columns}[T]
    \begin{column}{0.55\textwidth}
      \begin{minted}[fontsize=\scriptsize]{latex}
...
Iz izreka~\ref{izr:pitagora} izpeljemo
razdaljo med točkama v prostoru~[1].

\section*{Literatura}
[1] I. Vidav, \emph{Višja matematika I}, 2008.
      \end{minted}
    \end{column}
    \begin{column}{0.4\textwidth}
      \footnotesize
      Iz izreka 3.1 izpeljemo razdaljo med točkama v prostoru~[1].
      \razdelek{}{Literatura}
      [1] I. Vidav, \emph{Višja matematika I}, 2008.
    \end{column}
  \end{columns}
\end{frame}

\begin{frame}[fragile]{}
  \begin{columns}[T]
    \begin{column}{0.55\textwidth}
      \begin{minted}[fontsize=\scriptsize,highlightlines={1-2,8}]{latex}
Dokaz izreka je že v Evklidovih
Elementih~[1].
...
Iz izreka~\ref{izr:pitagora} izpeljemo
razdaljo med točkama v prostoru~[2].

\section*{Literatura}
[1] Evklid, \emph{Elementi}.

[2] I. Vidav, \emph{Višja matematika I}, 2008.
      \end{minted}
    \end{column}
    \begin{column}{0.4\textwidth}
      \footnotesize
      Dokaz izreka je že v Evklidovih Elementih~[1].

      \dots

      Iz izreka 3.1 izpeljemo razdaljo med točkama v prostoru~\alert{[1]}.
      \razdelek{}{Literatura}
      [1] Evklid, \emph{Elementi}.

      [2] I. Vidav, \emph{Višja matematika I}, 2008.
    \end{column}
  \end{columns}
\end{frame}

\begin{frame}[fragile]{}
  \begin{columns}[T]
    \begin{column}{0.55\textwidth}
      \begin{minted}[fontsize=\scriptsize,highlightlines={1-2},stripnl=false]{latex}


...
Iz izreka~\ref{izr:pitagora} izpeljemo
razdaljo med točkama v prostoru~[1].

\section*{Literatura}
[1] Evklid, \emph{Elementi}.

[2] I. Vidav, \emph{Višja matematika I}, 2008.
      \end{minted}
    \end{column}
    \begin{column}{0.4\textwidth}
      \footnotesize
      \invisible{Dokaz izreka je že v Evklidovih Elementih~[1].}

      \dots

      Iz izreka 3.1 izpeljemo razdaljo med točkama v prostoru~[1].
      \razdelek{}{Literatura}
      \alert{[1] Evklid, \emph{Elementi}.}

      [2] I. Vidav, \emph{Višja matematika I}, 2008.
    \end{column}
  \end{columns}
\end{frame}

\begin{frame}[fragile]{\texttt{literatura.bib}}
  \begin{columns}[T]
    \begin{column}{0.48\textwidth}
      \begin{minted}[fontsize=\scriptsize]{bibtex}
@book{evklid56,
  author    = {Evklid},
  title     = {The Thirteen Books
               of the Elements},
  note      = {prevedel T.~L.~Heath},
  publisher = {Dover},
  year      = 1956
}
      \end{minted}
    \end{column}
    \begin{column}{0.48\textwidth}
      \begin{minted}[fontsize=\scriptsize]{bibtex}
@book{vidav08,
  author    = {Ivan Vidav},
  title     = {Višja matematika I
               (12.~ponatis)},
  publisher = {DMFA--Založništvo},
  year      = 2008
}
@book{rudin87,
  author    = {Walter Rudin},
  title     = {Real and Complex Analysis},
  publisher = {McGraw-Hill},
  year      = 1987
}
      \end{minted}
    \end{column}
  \end{columns}
\end{frame}

\begin{frame}[fragile]{}
  \begin{minted}[fontsize=\footnotesize]{latex}
Dokaz izreka je že v Evklidovih Elementih~\cite{evklid56}.
...
Iz izreka~\ref{izr:pitagora} izpeljemo razdaljo med točkama
v prostoru~\cite{vidav08}.

\bibliographystyle{plain}
\bibliography{literatura}
  \end{minted}
\end{frame}

\begin{frame}[fragile]{}
  \begin{columns}[T]
    \begin{column}{0.55\textwidth}
      \begin{minted}{latex}
Iz izreka~\ref{izr:pitagora}
izpeljemo razdaljo med
točkama v prostoru~\cite{vidav}.
      \end{minted}
    \end{column}
    \begin{column}{0.4\textwidth}
    Iz izreka 3.1 izpeljemo razdaljo med točkama v prostoru~\textbf{[?]}.
    \end{column}
  \end{columns}
  \bigskip
  \begin{minted}[fontsize=\footnotesize]{text}
LaTeX Warning: Citation `vidav' on page 2 undefined on input line 31.
  \end{minted}
\end{frame}

\begin{frame}[fragile]{}
  \begin{minted}{bash}
$ pdflatex dokument.tex
$ bibtex dokument
$ pdflatex dokument.tex
$ pdflatex dokument.tex
  \end{minted}
\end{frame}

\begin{frame}[fragile]{}
  \begin{tabular}{lll}
    \mintinline{latex}|\bibliographystyle{plain}| & Elementih~[1] & prostoru~[2] \\
    \mintinline{latex}|\bibliographystyle{alpha}| & Elementih~[Evk56] & prostoru~[Vid08] \\
  \end{tabular}
  \bigskip
\end{frame}

\begin{frame}[fragile]{Vnosov ne tipkamo sami}
  \begin{center}
    \Large
    \href{https://zbmath.org/}{zbMATH Open \faLockOpen}

    \href{https://mathscinet.ams.org/}{MathSciNet \faLock}
  \end{center}
\end{frame}

\section{Lastni ukazi in okolja}

\begin{frame}[fragile]{Za pogoste izraze si definiramo ukaze}
  \begin{slaboenv}
    \begin{minted}{latex}
Kjer je matrika s koeficienti iz $\mathbb{R}^n$, $b \in \mathbb{R}^n$
je stolpec skalarjev, $x \in \mathbb{R}^n$ pa stolpec neznank.
    \end{minted}
  \end{slaboenv}

  \begin{dobroenv}
    \begin{minted}{latex}
\newcommand{\Rn}{\mathbb{R}^n}
Kjer je matrika s koeficienti iz $\Rn$, $b \in \Rn$ je stolpec
skalarjev, $x \in \Rn$ pa stolpec neznank.
    \end{minted}
  \end{dobroenv}
\end{frame}

\begin{frame}[fragile]{Matematične ukaze definiramo brez \texttt{\$}}
  \begin{slaboenv}
    \begin{minted}{latex}
\newcommand{\Rn}{$\mathbb{R}^n$}
Kjer je matrika s koeficienti iz \Rn, $b \in$ \Rn je stolpec
skalarjev, $x \in$ \Rn pa stolpec neznank.
    \end{minted}
    \newcommand{\Rn}{$\mathbb{R}^n$}
    Kjer je matrika s koeficienti iz \Rn, $b \in$ \Rn\ je stolpec skalarjev,
    $x \in$ \Rn\ pa stolpec neznank.
  \end{slaboenv}

  \begin{dobroenv}
    \begin{minted}{latex}
\newcommand{\Rn}{\mathbb{R}^n}
Kjer je matrika s koeficienti iz $\Rn$, $b \in \Rn$ je stolpec
skalarjev, $x \in \Rn$ pa stolpec neznank.
    \end{minted}
    \newcommand{\Rn}{\mathbb{R}^n}
    Kjer je matrika s koeficienti iz $\Rn$, $b \in \Rn$ je stolpec skalarjev,
    $x \in \Rn$ pa stolpec neznank.
  \end{dobroenv}
\end{frame}

\end{document}
